% ────────────────────────────────────────────
%  Masterprüfung — Präsentation (TU Graz)
%  Literatur kommt aus derselben Bibliothek wie die Masterarbeit.
% ────────────────────────────────────────────
\documentclass[aspectratio=169, 11pt]{beamer}
\usetheme{TUGraz}

\usepackage[utf8]{inputenc}
\usepackage{csquotes}
\usepackage{booktabs}
\usepackage{amsmath}
\usepackage{pgfplots}\pgfplotsset{compat=1.18}
% Literatur: biblatex + biber (natbib-Befehle wie \citep/\citet funktionieren weiter)
\usepackage[backend=biber, style=authoryear, natbib=true, maxcitenames=2, uniquename=false]{biblatex}
\addbibresource{references.bib}

\title{Titel der Masterarbeit}
\subtitle{Masterprüfung}
\author{Norbert Winter}
\institute{Institut für … \textbullet{} Technische Universität Graz}
\date{\today}

\begin{document}

\begin{frame}[plain]
  \titlepage
\end{frame}

\begin{frame}{Agenda}
  \tableofcontents
\end{frame}

\section{Motivation}

\begin{frame}{Warum dieses Thema?}
  \begin{itemize}
    \item Ausgangslage und Problem in einem Satz
    \item Relevanz für Forschung und Praxis
    \item Lücke in der bestehenden Literatur \citep{knuth1984texbook}
  \end{itemize}
  \vspace{4mm}
  \begin{block}{Forschungsfrage}
    Wie lässt sich \dots{} verbessern?
  \end{block}
\end{frame}

\section{Methodik}

\begin{frame}{Vorgehen}
  \begin{columns}[T]
    \begin{column}{0.48\textwidth}
      \begin{enumerate}
        \item Literaturrecherche
        \item Konzept \& Modell
        \item Implementierung
        \item Evaluation
      \end{enumerate}
    \end{column}
    \begin{column}{0.48\textwidth}
      \begin{exampleblock}{Zielfunktion}
        \[ \mathcal{L}(\theta) = -\frac{1}{N}\sum_{i=1}^{N}\log p_\theta(y_i\mid x_i) \]
      \end{exampleblock}
    \end{column}
  \end{columns}
\end{frame}

\section{Ergebnisse}

\begin{frame}{Ergebnisse im Überblick}
  \centering
  \begin{tikzpicture}
    \begin{axis}[width=0.78\textwidth, height=5.4cm, xlabel={Epoche}, ylabel={Genauigkeit},
      grid=major, grid style={black!8}, legend pos=south east, legend style={draw=none, font=\small},
      every axis plot/.append style={very thick}]
      \addplot[color=tured, mark=*] coordinates {(1,0.52) (2,0.67) (3,0.75) (4,0.81) (5,0.84) (6,0.86)};
      \addplot[color=tugray, mark=square*] coordinates {(1,0.48) (2,0.58) (3,0.66) (4,0.71) (5,0.74) (6,0.75)};
      \legend{Unser Ansatz, Baseline}
    \end{axis}
  \end{tikzpicture}
\end{frame}

\begin{frame}{Kennzahlen}
  \centering
  \begin{tabular}{lrr}
    \toprule
    Methode     & Genauigkeit & Laufzeit \\
    \midrule
    Baseline    & 75{,}1\,\%  & 12{,}4\,h \\
    \alert{Unser Ansatz} & \alert{86{,}3\,\%} & \alert{9{,}8\,h} \\
    \bottomrule
  \end{tabular}
\end{frame}

\section{Fazit}

\begin{frame}{Fazit \& Ausblick}
  \begin{itemize}
    \item<1-> Kernaussage 1
    \item<2-> Kernaussage 2
    \item<3-> Offene Fragen und nächste Schritte
  \end{itemize}
\end{frame}

\danke

% ── Backup / Anhang ──
\begin{frame}{Literatur}
  \scriptsize
  \printbibliography[heading=none]
\end{frame}

\end{document}
